\documentclass[12pt, a4paper, oneside]{ctexart}

\usepackage{amsmath, amsthm, amssymb, appendix, bm, fancyhdr, graphicx, geometry, mathrsfs, zhnumber}
\usepackage[bookmarks=true, colorlinks, citecolor=black, linkcolor=black]{hyperref}  %超链接
\usepackage[framed, numbered, autolinebreaks, useliterate]{mcode}  %matlab代码
\linespread{1.2}    %行距
\geometry{left=2.5cm, right=2.5cm, top=2.5cm, bottom=2.5cm} %边距
\graphicspath{{./graphics/}}    %图片文件夹

\renewcommand{\thesection}{\zhnum{section}} %一级标题标号

\renewcommand{\thesubsection}{\arabic{section}.\arabic{subsection}}
\renewcommand{\theequation}{\arabic{section}.\arabic{equation}}
\renewcommand{\thetable}{\arabic{section}.\arabic{table}}
\renewcommand{\thefigure}{\arabic{section}.\arabic{figure}}

\begin{document}

    \pagestyle{empty}
    \setcounter{page}{0}

    \begin{center}
        \Large{\textbf{\LaTeX 建模论文模板3.0}}
    \end{center}

    \begin{center}
        \Large{\textbf{摘要}}
    \end{center}

    这里是摘要。

    \textbf{关键词：}这里是关键词；这里是关键词。

    \newpage
    \tableofcontents
    \newpage
    \setcounter{page}{1}
    \pagenumbering{arabic}
    \pagestyle{plain}
    \fancyfoot[C]{\thepage}

    \section{问题背景与重述}

    \subsection{问题背景}

    编写该模板是为了应对接下来所面对的一切中文数学建模竞赛。

    众所周知，Word的公式排版奇丑无比，并且小节和公式的编号较为麻烦。
    相比之下，我更喜欢\LaTeX 干净简洁的风格。
    为了彻底摆脱Word，特地编写了该模板，用于接下来的各种竞赛。

    \subsection{问题重述}

    我们所需要解决的问题如下。
    \begin{itemize}
        \item 制作出一个适用于中文建模竞赛的\LaTeX 模板；
        \item 在模板中，应当能够使用表格、图片、公式等对象。
    \end{itemize}

    \section{问题分析}

    \subsection{问题1的分析}

    在这里写问题1的分析。

    \subsection{问题2的分析}

    在这里写问题2的分析。

    \section{模型准备}

    \subsection{模型假设}

    为了建立模型，我们提出如下的假设。

    \begin{enumerate}
        \item 这里是第一条假设。
        \\\textbf{理由：}这里是作出第一条假设的理由。
        \item 这里是第二条假设。
        \\\textbf{理由：}这里是作出第二条假设的理由。
        \item 这里是第三条假设。
        \\\textbf{理由：}这里是作出第三条假设的理由。
        \item 这里是第四条假设。
        \\\textbf{理由：}这里是作出第四条假设的理由。
    \end{enumerate}

    \subsection{符号说明}

    所使用的符号及说明如表\ref{table1}所示。

    \begin{table}[h]
        \caption{符号说明}\label{table1}
        \centering
        \begin{tabular}{clc}
            \hline
            \textbf{符号} & \textbf{说明}        & \textbf{单位}   \\ \hline
            符号1         & 这里是符号1的说明。  & 单位             \\
            符号2         & 这里是符号2的说明。  & 单位             \\
            符号3         & 这里是符号3的说明。  & 单位             \\ \hline
        \end{tabular}
    \end{table}

    \section{模型的建立与求解}

    \subsection{模型1}

    针对问题1，建立了模型1。

    其中，公式的书写方式如下。
    \begin{equation}
        \label{eq1}
        {\rm{e}}^{i\theta}=\cos\theta+i\sin\theta.
    \end{equation}
    公式\ref{eq1}就是大名鼎鼎的Euler公式。

    \subsection{模型2}

    针对问题2，建立了模型2。

    在论文中可能需要插入图片，在这里插入图片的方式如下。

    \begin{figure}[htbp]
        \centering
        \includegraphics[width=10cm]{灵梦.jpeg}
        \caption{微生物}\label{fig1}
    \end{figure}

    图\ref{fig1}是在实验室中，科学家拿着微生物的照片。

    \section{结果的分析与检验}

    \subsection{问题的结果}

    在这里写问题的结果。

    \subsection{模型的检验}

    在这里写对模型的检验。

    \section{模型的优缺点分析}

    \subsection{模型的优点}

    该模型具有如下的优点。
    \begin{itemize}
        \item 优点1；
        \item 优点2；
        \item 优点3。
    \end{itemize}

    \subsection{模型的缺点与改进}

    与此同时，该模型也具有如下的缺点。
    \begin{itemize}
        \item 缺点1；
        \item 缺点2。
    \end{itemize}
    同时，在这里给出进一步优化模型的思路。

    \begin{thebibliography}{99}
        \bibitem{a}作者. \emph{文献}[M]. 地点:出版社,年份.
        \bibitem{b}作者. \emph{文献}[M]. 地点:出版社,年份.
    \end{thebibliography}

    \begin{center}
        \Large{\textbf{附录}}
    \end{center}

    \begin{appendices}
        \renewcommand{\thesection}{\Alph{section}}
        \section{所用软件}
            论文使用\LaTeX 排版。
        \section{代码}
            所使用的代码如下。
    \begin{lstlisting}
        Hello. 
    \end{lstlisting}
    \end{appendices}

\end{document}
